\subsection{Changing the phase and restricting the group}\label{sec:blocktransport}

Two choices enter a decomposition indexed by Euler values. The phase orders the distinct values. The group can force several labels to have one value. Changing the phase exchanges entire blocks, whereas passing to a subgroup can separate the labels inside a block.

\begin{proposition}\label{prop:blocktransport}
Fix an action of a finite group $G$ by automorphisms on a smooth complex Fano variety and a fixed semisimple parameter on a compatible calibrated branch. Suppose that, for one lifted admissible phase, a full exceptional collection realizes the individual Gamma columns in phase order. Assume the group permutes its objects according to the idempotent labels and that distinct objects of equal Euler value are mutually orthogonal in every degree. Then such a collection exists for every admissible lifted phase at the same parameter. More generally, the collection continues along a chosen path of $G$-fixed semisimple parameters and lifted phases on the same calibrated branch, starting at the marked pair and ending at an admissible phase, provided the partition into equal values is locally constant. At a nearby generic point of the fixed germ for a subgroup, each old equal-value block refines into its subgroup orbits with the same objects.
\end{proposition}

\begin{proof}
Use local canonical coordinates $u_s$ labelled by the idempotents. On the fixed germ $u_s=u_{gs}$ for every $g$. Imposing the additional equalities of the chosen partition leaves one independent complex coordinate $v_B$ for each block $B$. Distinct block coordinates are unequal. Along a path with lifted phase $\varphi$, the wall for two blocks is consequently
\[
             \operatorname{Im}\bigl(e^{-i\varphi}(v_B-v_C)\bigr)=0 .
\]
The real differentials for two different unordered pairs are independent in the transverse height coordinates. After a small deformation fixing admissible endpoints, a compact path therefore has finitely many transverse crossings, one pair of blocks at a time. For a change of phase at the original parameter we may make this deformation within a small coordinate ball and return to the original parameter at the last phase. In every chamber the normalized columns, with their signs transported, are flat by \cite[Proposition~2.5.1]{GGI}.

Consider a crossing of two distinct values carrying label sets $B$ and $C$. To determine the individual columns, resolve each value in its small coordinate disk, temporarily separating the labels in the height direction. Tiny offsets, constant in the rotated coordinate, may be chosen so every difference of an offset in $B$ and one in $C$ is distinct. On a short crossing interval the height difference of the original values is strictly monotone. Thus each pair from the two blocks crosses once. Their horizontal differences all have the same sign, since the two distinct block values lie on the same oriented line at the wall. Suppose $B$ initially precedes $C$ in decreasing imaginary height. With $x_D=\operatorname{Re}(e^{-i\varphi}v_D)$ at the wall, $x_C>x_B$ right-mutates $B$ through $C$, and $x_C<x_B$ left-mutates $C$ through $B$. Thus all the elementary exchanges have the same direction. The ordered-cover description of the marked system and the extension of the normalized columns back to the tied endpoints give the change on the original individual labels \cite[Propositions~2.5.1, 4.3.4 and 4.3.8]{GGI}.

For a right mutation let $\mathcal C$ be the admissible subcategory generated by the objects $E_c$ in the block that stays fixed. For each $b\in B$ the new object is determined by the coevaluation triangle
\[
 R_{\mathcal C} E_b\longrightarrow E_b\longrightarrow
   \bigoplus_{c\in C}\operatorname{RHom}(E_b,E_c)^\vee\otimes E_c
   \longrightarrow (R_{\mathcal C}E_b)[1].
\]
The summands give the left adjoint projection to $\mathcal C$ because the $E_c$ are mutually orthogonal exceptional objects. Successive right mutations past the individual $E_c$ give this same fibre: projecting onto one summand leaves every other summand unchanged. Taking classes in the triangle gives exactly the successive Euler mutations of the analytic columns. Right mutation restricts to an equivalence between the two orthogonals to $\mathcal C$, taking the old objects in $\mathcal C^\perp$ to its left orthogonal, so it preserves all morphism complexes within the moving block. The inverse crossing uses the left mutation triangle. Each exchange preserves generation. Since $\mathcal C$ is stable under the group, its inclusion and adjoints intertwine pullback; hence the mutation commutes with pullback up to isomorphism \cite[Lemma~3.24]{ESS}. The resulting objects are permuted with their labels. The signs relative to orbit-compatible normalized columns are constant on each orbit by equivariance of the framing. This proves the continuation assertion along the chosen path.

For the last assertion fix the phase at the original parameter and restrict to a sufficiently small coordinate ball. For a subgroup $J$, there is now one independent coordinate for each $J$-orbit. A generic point of its fixed germ separates these coordinates, with their heights distinct; comparisons between previously distinct values retain their order. The sectorial frame extends through the starting ties \cite[Proposition~2.5.1]{GGI}. Thus the original objects have the prescribed classes at the nearby parameter. The only reorderings needed are within old blocks, where all off-diagonal morphism complexes vanish.
\end{proof}

At a new coincidence reached from a stratum with separate values, the orthogonality hypothesis for the merged block requires its own geometric argument.
